\documentclass[12pt,a4paper]{article}

\usepackage[utf8]{inputenc}
\usepackage[T1]{fontenc}
\usepackage[ngerman]{babel}
\usepackage{url}

\begin{document}
	
	\begin{titlepage}
		\centering
		
		{\Large Hochschule für Technik und Wirtschaft Berlin\par}
		\vspace{1cm}
		
		{\Large Lehrveranstaltung: Softwaretechnik\par}
		\vspace{2cm}
		
		{\Huge\bfseries Inventory Management Systems\par}
		\vspace{0.5cm}
		
		{\Large Meilenstein MS01\par}
		\vspace{2cm}
		
		\begin{flushleft}
			\textbf{Projekt:} Inventory Management System / WareHouse \\[0.3cm]
			\textbf{Quelle:} \url{https://github.com/ZeroCoolHacker/Inventory} \\[0.3cm]
			\textbf {Authors:}\\ [0.3cm]
			\textbf Rey Krisna Tri Dalmesta \\[0.3cm]
			\textbf Rey Indra Tri Dalmesta (s0587277) \\[0.3cm]
			\textbf SoftwareTechnik \\[0.3cm]
		\end{flushleft}
		
		\vfill
		
	\end{titlepage}
	
	\tableofcontents
	\newpage
	
	\section{Informelle Beschreibung der Software}
	
	Die untersuchte Software ist eine Lagerverwaltungssoftware mit dem Namen
	\texttt{WareHouse}. Das Projekt ist eine Desktop-Anwendung, die in C++ mit dem
	Qt-Framework entwickelt wurde. Die Anwendung verwendet eine lokale
	SQLite-Datenbank zur Speicherung der Daten.
	
	\subsection{Zweck der Software}
	
	Der Zweck der Software ist die Verwaltung einfacher Lagerprozesse. Mit der
	Anwendung können Artikel, Lieferanten, Einkaufsvorgänge, Verkaufsvorgänge und
	Berichte verwaltet werden.
	
	Die Software unterstützt folgende zentrale Funktionen:
	
	\begin{itemize}
		\item Verwaltung von Artikeln
		\item Verwaltung von Lieferanten
		\item Erfassung von Einkäufen
		\item Erfassung von Verkäufen
		\item Aktualisierung des Lagerbestands nach Einkauf oder Verkauf
		\item Anzeige von Einkaufs- und Verkaufsberichten
		\item Speicherung der Daten in einer SQLite-Datenbank
	\end{itemize}
	
	Die Software eignet sich daher für ein kleines Lager oder ein einfaches
	Geschäft, in dem Warenbestand, Lieferanten und Rechnungen verwaltet werden
	müssen. Es handelt sich nicht um ein vollständiges ERP-System, sondern um eine
	überschaubare Anwendung für grundlegende Lagerverwaltung.
	
	\subsection{Nutzerrollen}
	
	Im Quellcode gibt es keine Benutzerverwaltung mit Login oder verschiedenen
	technischen Benutzerkonten. Für das fachliche Modell können jedoch folgende
	Nutzerrollen beschrieben werden:
	
	\begin{description}
		\item[Mitarbeiter]
		Der Mitarbeiter ist der Hauptnutzer der Software. Er bedient die grafische
		Oberfläche und führt die Lagerprozesse aus. Dazu gehören das Anlegen von
		Artikeln und Lieferanten, das Erfassen von Einkäufen und Verkäufen sowie
		das Anzeigen von Berichten.
		
		\item[Lieferant]
		Der Lieferant ist eine externe Rolle. Lieferantendaten werden im System
		gespeichert und bei Einkaufsvorgängen verwendet. Der Lieferant benutzt die
		Software jedoch nicht direkt.
		
		\item[Kunde]
		Der Kunde ist ebenfalls eine externe Rolle. Bei Verkaufsvorgängen werden
		Kundendaten beziehungsweise Käuferinformationen verwendet. Auch der Kunde
		arbeitet nicht direkt mit der Software, sondern wird durch den Mitarbeiter
		im System erfasst.
		
		\item[System]
		Das System speichert die eingegebenen Daten, prüft Eingaben, aktualisiert
		den Lagerbestand und liest Daten aus der SQLite-Datenbank für Berichte.
	\end{description}
	
	\subsection{Quelle der Software}
	
	Die Software basiert auf einem bestehenden Open-Source-Projekt auf GitHub:
	
	\begin{center}
		\url{https://github.com/ZeroCoolHacker/Inventory}
	\end{center}
	
	Für die Analyse wurde der Quellcode des Projekts untersucht. Wichtige
	Bestandteile sind die Qt-Klassen für das Hauptfenster, Artikel, Lieferanten,
	Einkäufe, Verkäufe und Berichte sowie die lokale SQLite-Datenbank.
	

	\section{Geplante Erweiterungen und Verbesserungen}
	
	Für die zweite Projektphase ist geplant, die bestehende Lagerverwaltungssoftware
	gezielt zu erweitern. 
	
	\subsection{Kundenverwaltung}
	
	Die Software soll um eine Kundenverwaltung erweitert werden. Bisher stehen vor
	allem Artikel, Lieferanten, Einkäufe und Verkäufe im Mittelpunkt. Durch eine
	eigene Kundenverwaltung können Kundendaten gespeichert und mit Verkaufsvorgängen
	verbunden werden.
	
	Geplante Funktionen:
	
	\begin{itemize}
		\item Kunden anlegen
		\item Kundendaten bearbeiten
		\item Kunden löschen
		\item Kunden mit Verkäufen verknüpfen
		\item Verkaufshistorie eines Kunden anzeigen
	\end{itemize}
	
	\subsection{Bestandswarnung}
	
	Das System soll eine Warnung anzeigen, wenn der Lagerbestand eines Artikels
	unter einen festgelegten Mindestwert fällt. Dadurch kann der Nutzer rechtzeitig
	erkennen, welche Artikel nachbestellt werden müssen.
	
	Geplante Funktionen:
	
	\begin{itemize}
		\item Mindestbestand pro Artikel speichern
		\item niedrigen Lagerbestand automatisch erkennen
		\item Warnmeldung oder Markierung in der Artikelübersicht anzeigen
		\item betroffene Artikel in einem Bericht anzeigen
	\end{itemize}
	
	\subsection{Inventarhistorie}
	
	Zusätzlich soll eine Historie für Bestandsänderungen eingeführt werden. Jede
	Änderung am Lagerbestand soll nachvollziehbar gespeichert werden. Dadurch kann
	später geprüft werden, warum sich der Bestand eines Artikels verändert hat.
	
	Geplante Informationen in der Historie:
	
	\begin{itemize}
		\item betroffener Artikel
		\item Art der Änderung, zum Beispiel Einkauf oder Verkauf
		\item geänderte Menge
		\item Bestand vor und nach der Änderung
		\item Zeitpunkt der Änderung
	\end{itemize}
	
	\subsection{Lieferantenbestellungen}
	
	Die Software soll außerdem um einfache Lieferantenbestellungen erweitert werden.
	Wenn ein Artikel einen niedrigen Bestand hat, soll daraus eine Bestellung bei
	einem Lieferanten erstellt werden können.
	
	Geplante Funktionen:
	
	\begin{itemize}
		\item Bestellung für einen Lieferanten erstellen
		\item mehrere Artikel zu einer Bestellung hinzufügen
		\item offene Bestellungen anzeigen
		\item Bestellung als erhalten markieren
		\item Lagerbestand nach Wareneingang aktualisieren
	\end{itemize}
	
	\subsection{Benutzerverwaltung und Rollen}
	
	Die bestehende Software besitzt keine technische Benutzerverwaltung. Zwar können
	im Use-Case-Modell mehrere fachliche Akteure wie Mitarbeiter, Lieferant und
	Kunde unterschieden werden, diese sind jedoch keine implementierten
	Benutzerkonten mit eigenen Zugriffsrechten. Eine geplante Erweiterung ist daher
	die Einführung einer Benutzerverwaltung mit Login-Funktion und Rollen wie Admin
	und Mitarbeiter.
	
	Geplante Rollen:
	
	\begin{description}
		\item[Admin]
		Der Admin kann Benutzer verwalten, Artikel und Lieferanten bearbeiten,
		Berichte ansehen und kritische Daten löschen.
		
		\item[Mitarbeiter]
		Der Mitarbeiter kann Artikel ansehen, Einkäufe und Verkäufe erfassen und
		Berichte abrufen. Kritische Verwaltungsfunktionen sind eingeschränkt.
	\end{description}
	
	Geplante Funktionen:
	
	\begin{itemize}
		\item Login
		\item Logout
		\item Benutzer anlegen
		\item Benutzer bearbeiten
		\item Benutzer löschen
		\item Rechte je nach Rolle einschränken
	\end{itemize}
	
	
	
\end{document}
